% ============================================================================
% ABSTRACT — Journal of Big Data (expanded, max 350 words)
% ============================================================================
\begin{abstract}
\noindent A data engineer maintaining an enterprise warehouse encounters
structural questions every day: which upstream tables break if this column
changes, which tables share a lineage chain, whether two tables are even
connected.  These questions are not about SQL syntax.  They are about
graph topology, and answering them correctly requires traversing
foreign-key paths, detecting connected components, and tracing
multi-hop lineage chains across heterogeneous schema graphs.

We ask a direct question: can frontier LLMs do this?  To find out, we
built \textbf{DW-Bench}, a benchmark of 1{,}479 questions over five real
and synthetic data warehouse schemas, spanning schema-structure reasoning
and row-level provenance tracing across 14 question subtypes.

We evaluated six baselines, from static schema injection to interactive
tool-use, against three frontier models (Gemini~2.5 Flash, DeepSeek-V3,
Qwen2.5-72B).
Tool-augmented baselines reach 89--90\% accuracy on schema-level
questions, outperforming static methods by 13--14 percentage points.
But static baselines exhibit a \textit{Triviality Illusion}: easy
lookup subtypes dominate the question pool and inflate aggregate scores,
masking systematic failures on compositional reasoning.
Hard compositional tasks plateau at 60\% even with full tool access,
while the Oracle clears 99\%.
The hardest subtype, \texttt{combined\_impact} (requiring lineage
traversal composed with FK propagation across multiple hops), reaches
only 12\% exact match with the best tool-augmented baseline.
Budget ablation confirms the bottleneck is a reasoning deficit,
not a tool or context-length limitation.

A controlled obfuscation experiment reveals what is happening:
tool-augmented baselines are robust to name randomization
(accuracy drops below 4\%), while static text-injection baselines
lose 9--32\% when table names are anonymized.
One group is reasoning about structure; the other is pattern-matching
against schema names.

All datasets, evaluation code, question generators, and baseline
implementations are publicly available at
\url{https://github.com/AJamal27891/dw-bench}, providing a reproducible
foundation for measuring progress on structural reasoning in
data-intensive systems.
\end{abstract}
